\section{Ba$_5$Y$_{12}$Zn[O(SiO$_4$)]$_8$ 的类似物判读}

本次检索没有在 OpenAlex、Crossref 或可访问的全文语料中找到 Ba$_5$Y$_{12}$Zn[O(SiO$_4$)]$_8$ 的原始晶体结构条目。检索到的“钇钡硅酸盐”主要是发光薄膜和氧基磷灰石类材料，化学计量、空间群和 Zn 占位均不等同于目标式 \cite{gupta2021combustion,yang2021cyan}。因此，当前最稳妥的表述是“待验证的结构类似物候选”，而不是“已知同构物”。

仍可提出分层比较。BaZn$_2$Si$_2$O$_7$ 提供 ZnO$_4$—SiO$_4$ 角共享和 Ba 多面体的直接参照 \cite{lin1999phase,available2020materials,kaiser2002crystal}；钇钡硅酸盐氧基磷灰石文献提供 Y/Ba 与孤立硅酸根共存的化学邻域，但并未证明相同拓扑 \cite{gupta2021combustion,yang2021cyan,wei2013a}。若未来获得目标相粉末或单晶数据，应依次核对硅酸根聚合度、Zn 配位数、Y 多面体连接、氧空位位置和空间群，而不是只比较化学式中元素比例。

Ba$_5$Y$_{12}$Zn[O(SiO$_4$)]$_8$ 的高 Y 含量还提示明显的电荷补偿问题。形式价态计算要求额外 O$^{2-}$ 或阳离子混合价态；任何合成尝试都必须用热重、XPS\slash XANES 或中子衍射确认氧计量和价态。BaZn$_2$Si$_2$O$_7$ 中观察到的热致相变可作为热处理窗口的启发，但不能直接外推到 Y-rich 体系。

\begin{table}[t]
\centering
\caption{结构类似物比较的证据层级。}
\begin{tabular}{P{0.22\linewidth}P{0.31\linewidth}P{0.30\linewidth}}
\toprule
比较层级 & BaZn$_2$Si$_2$O$_7$ 的依据 & 对 Ba$_5$Y$_{12}$Zn[O(SiO$_4$)]$_8$ 的含义 \\
\midrule
硅酸根聚合度 & ZnO$_4$/SiO$_4$ 角共享网络 \cite{lin1999phase} & 可作为待核对的拓扑起点，非同构证明 \\
阳离子配位 & Ba 五配位、Zn 四配位 \cite{available2020materials} & 需实测 Y 和 Zn 配位，不能按离子半径指定 \\
热稳定性 & 约 280--300~$^\circ$C 相变 \cite{thieme2015ba1} & 仅能提供热分析程序参照 \\
\bottomrule
\end{tabular}
\end{table}
